%%=============================================================================
%% De inhoud van je projectvoorstel
%%=============================================================================
%%
%% Zeven rubrieken, samen hoogstens twee bladzijden. Houd het kort: per
%% rubriek volstaan één tot drie zinnen. Een voorstel is geen rapport - het
%% moet je begeleider in vijf minuten laten beslissen of dit een geschikt
%% graduaatsproject is.
%%
%% Stem je voorstel af met je mentor VOORDAT je het indient. Je mentor
%% bevestigt dat dit is wat het bedrijf bedoelt en dat de tien werkdagen
%% haalbaar zijn.

%---------- 1. Context en probleemstelling ------------------------------------

\section{Context en probleemstelling}%
\label{vst:context}

Bij Endare moeten medewerkers regelmatig gevoelige gegevens, zoals
wachtwoorden, API-sleutels en toegangsgegevens, met klanten delen. Vandaag
gebeurt dat via e-mail of chat. Daar blijven die gegevens onbeperkt staan en
zijn ze leesbaar voor iedereen die later toegang krijgt tot die berichten.
Externe diensten zoals onetimesecret.com zijn een alternatief, maar dan ligt
de data bij een derde partij waar Endare geen controle over heeft.

%---------- 2. Doelstelling ---------------------------------------------------

\section{Doelstelling}%
\label{vst:doelstelling}

Op het einde van dit project bestaat er een OneTimeSecret-platform dat Endare
zelf host en beheert. Endare-medewerkers melden aan met hun Google
Workspace-account en maken een secret aan met inhoud, levensduur en
wachtwoord. Daarbij krijgen ze een unieke link voor de klant. Het project is
geslaagd wanneer aan de volgende voorwaarden voldaan is:
\begin{itemize}
  \item een klant kan het secret via die link en het wachtwoord exact één
        keer bekijken;
  \item het secret is daarna aantoonbaar van de server verwijderd;
  \item de inhoud staat op geen enkel moment onversleuteld in de databank.
\end{itemize}

%---------- 3. Technologieën --------------------------------------------------

\section{Technologieën}%
\label{vst:technologieen}

Ik bouw een webapplicatie met een TypeScript/Node.js-backend, een
React-frontend en een PostgreSQL-databank. De applicatie draait in Docker
zodat ze op de infrastructuur van Endare gehost kan worden.

Voor de authenticatie gebruik ik Google OAuth 2.0 / OpenID Connect, beperkt
tot het Endare-domein. Voor de encryptie leid ik een sleutel af van het
wachtwoord met Argon2id en versleutel ik de inhoud met AES-256-GCM, telkens
met bewezen cryptografische bibliotheken. Als uitbreiding koppel ik het
platform aan de API van Passbolt.

Nieuw voor mij zijn het toepassen van cryptografie in een echte applicatie,
OpenID Connect en de Passbolt API. Ik zet AI in als hulp bij het opzoeken
van best practices, het nalezen van code en het schrijven van tests. Er zit
geen AI in de applicatie zelf.

De applicatie verwerkt persoonsgegevens: naam en e-mailadres van
medewerkers. De secrets kunnen bovendien gevoelige klantgegevens bevatten,
dus de AVG speelt mee.

%---------- 4. Afbakening -----------------------------------------------------

\section{Afbakening}%
\label{vst:afbakening}

Mijn dagelijkse stagewerk bij Endare bestaat uit andere opdrachten, zoals de
integratie tussen ELLA en Odoo. Mijn graduaatsproject is het afgebakende
geheel rond het OneTimeSecret-platform. Ik analyseer zelfstandig de vereisten
en de beveiligingsrisico's, en ik ontwerp:
\begin{itemize}
  \item de architectuur;
  \item het datamodel;
  \item het encryptieschema, met de keuze en verantwoording van het
        sleutelafleidings- en encryptie-algoritme;
  \item de flow voor eenmalig bekijken en verwijderen.
\end{itemize}
Daarna realiseer, test en deploy ik het platform. De Passbolt-integratie is
een uitbreiding die ik enkel aanpak als de kern tijdig af is.

Buiten de afbakening vallen:
\begin{itemize}
  \item accounts of een login voor klanten;
  \item het delen van bestanden of bijlagen;
  \item uitgebreid rollen- en rechtenbeheer;
  \item een beheerdashboard met statistieken;
  \item integraties met andere tools dan Passbolt.
\end{itemize}
%---------- 5. Verwachte meerwaarde -------------------------------------------

\section{Verwachte meerwaarde}%
\label{vst:meerwaarde}

Endare kan gevoelige gegevens veilig en traceerbaar met klanten delen, zonder
dat die in mailboxen of chatgeschiedenis achterblijven en zonder afhankelijk
te zijn van een externe dienst. Na mijn stage blijft er een gedocumenteerd,
gedeployd intern platform over dat het team dagelijks kan gebruiken. Het
vormt ook een basis voor verdere uitbreidingen, zoals de koppeling met
Passbolt.

%---------- 6. Planning en werkbelasting --------------------------------------

\section{Planning en werkbelasting}%
\label{vst:planning}

\begin{table}[htbp]
  \centering
  \small
  \begin{tabular}{@{}p{0.3\columnwidth}p{0.47\columnwidth}r@{}}
    \toprule
    \textbf{Fase} & \textbf{Wat lever je op} & \textbf{Uren} \\
    \midrule
    Analyse en vereisten         & vereisten, dreigingsanalyse                  & 10  \\
    Ontwerp                      & architectuur, datamodel, encryptieschema     & 10  \\
    Secrets aanmaken             & versleutelde opslag, unieke link             & 20  \\
    Eenmalig bekijken            & ontsleuteling, verwijdering, verloop         & 15  \\
    Deployment                   & Docker-setup op Endare-infrastructuur        & 10  \\
    Testen en validatie          & geautomatiseerde tests, securitycontrole     & 10  \\
    Rapport, poster en demo      & de op te leveren documenten                  & 10  \\
    Uitbreiding Passbolt         & placeholders, ophalen bij bekijken           & 15  \\
    \midrule
    \textbf{Totaal}              &                                              & \textbf{100} \\
    \bottomrule
  \end{tabular}
  \caption[Planning en werkbelasting]{\label{tab:vst-planning}%
    Werkuren.}
\end{table}

%---------- 7. Risico's -------------------------------------------------------

\section{Risico's}%
\label{vst:risicos}

\begin{itemize}
  \item \textbf{Fouten in de encryptie maken het platform onveilig.} Ik
        gebruik daarom enkel vertrouwde bibliotheken, geen zelfgeschreven
        cryptografie, en laat het encryptieschema nalezen door een ervaren
        collega bij Endare.
  \item \textbf{Aanmaken project en JIRA bord.} Ik z.s.m al aan om het project op bitbucket en JIRA te zetten.
  \item \textbf{De kern neemt meer tijd dan gepland.} Passbolt integreren aan de kat schuiven en als eventuele uitbreiding gebruiken.
\end{itemize}